There are $15$ diagonals in one direction, and at most one bishop may occupy each. The two extreme diagonals each consist of one corner square; those two corners lie on the same diagonal in the other direction, so both cannot be occupied. Therefore at most $14$ bishops are possible.

Place bishops on all eight squares of the first row and on the six noncorner squares of the last row. Within a row no bishops attack. Between rows the vertical separation is $7$ while the horizontal separation is at most $6$, so there is no diagonal attack. Thus the maximum is $\boxed{14}$.
\begin{asy}
texpreamble("\usepackage{amsmath,amssymb}");
size(240);
defaultpen(fontsize(9pt));
for(int x=0;x<8;++x)for(int y=0;y<8;++y){filldraw(box((x,y),(x+1,y+1)),(x+y)%2==0?white:lightgray);}for(int x=0;x<8;++x)label("B",(x+.5,.5));for(int x=1;x<7;++x)label("B",(x+.5,7.5));
\end{asy}
